% Unit 26 — early half-angle substitution t=tan(x/2)
\input{technique-note.tex}
\title{Integration Bee --- Unit 26:\\ Half-angle substitution}
\author{UK Uni Integration Bee prep\\ Techniques Guide \S1.6 T-sub}
\date{}
\begin{document}
\maketitle
\section*{Status}
Tags: \texttt{yang:4.5}; techniques: \texttt{weierstrass}.
Guide section 1.6 is this move. The syllabus lists it in section 1, so it is an early tool: use it as soon as the integrand is a rational function of sine and cosine. Yang \S4.5 is a later chapter under the name Weierstrass. It is the same substitution, not a new one. A first pass sits in unit 09. This file is the note a miss names.
\section{Official anchors}
\begin{itemize}
\item Techniques Guide \S1.6, the T-substitution $t=\tan(x/2)$.
\item Syllabus \S1 Weierstrass (\texttt{sources/UK\_University\_Integration\_Bee\_Syllabus.pdf}).
\item Yang Topics Guide 2nd Ed \S4.5, the same three replacements.
\end{itemize}
\section{Learning goals}
Turn a rational function of sine and cosine into a rational function of $t$.
The substitution is $t=\tan(x/2)$. Check the endpoints before you expand. Refuse the $[0,\pi]$ value on a shorter interval, and refuse a definite value when the blank wants an antiderivative. A plain sine power is a reduction, unit 20, not this substitution. A rational function of $\tan x$ with a $\sec^2 x$ already present is $t=\tan x$, which is unit 09's ordinary substitution.

\section{The three replacements}
The half-angle formulas give
\[
\sin x=\frac{2t}{1+t^2},\qquad
\cos x=\frac{1-t^2}{1+t^2},\qquad
dx=\frac{2\,dt}{1+t^2},
\]
where $t=\tan(x/2)$. They are the double-angle formulas written in $t$, together with $x=2\arctan t$. The quotient of the first two is
\[
\tan x=\frac{2t}{1-t^2}.
\]
Use the three replacements, and this fourth when a tangent is present, only when the integrand is rational in $\sin x$ and $\cos x$. The map $t=\tan(x/2)$ is continuous on $(-\pi,\pi)$. On $[0,\pi/2]$ the new limits are finite. If an odd multiple of $\pi$ lies in the interior, split there: the half-angle tangent is discontinuous. If the endpoint is $x=\pi$, the integral in $t$ is improper. Take the limit. Do not paste an antiderivative across that point.

\section{Worked --- \texttt{U26-D01}}
\[
\int_0^{\pi/2}\frac{1}{1+\sin x}\,dx.
\]
The limits $x=0$ and $x=\pi/2$ become $t=0$ and $t=1$.
\[
1+\sin x=\frac{(1+t)^2}{1+t^2},
\]
so the integrand collapses to $2/(1+t)^2$.
\[
\int_0^1\frac{2}{(1+t)^2}\,dt=\Bigl[\frac{-2}{1+t}\Bigr]_0^1=1.
\]
That is the stored answer. The worked dose stops there.

\section{Faded --- \texttt{U26-D08}}
\[
\int_0^{\pi/2}\frac{dx}{2+\sin x}.
\]
The blank is the pair of limits in $t$. They are $0$ and $1$, not $0$ and $\infty$. After that line is written, uncover the stored answer $\sqrt{3}\,\pi/9$.

\section{Bare --- \texttt{U26-D02}}
Set
\[
\int_0^{\pi}\frac{dx}{2+\cos x}
\]
and show nothing else. The bare dose hides $a$ and $b$. After a line is written, uncover $a=2$, $b=1$, and
\[
\int_0^{\pi}\frac{dx}{a+b\cos x}=\frac{\pi}{\sqrt{a^2-b^2}}
\]
for $a>|b|$. The stored answer is $\pi/\sqrt{3}$. The feed prints $\sqrt{3}\,\pi/3$.

As $x$ runs over $[0,\pi)$, $t=\tan(x/2)$ runs over $[0,\infty)$, and
\[
a+b\cos x=\frac{(a-b)t^2+(a+b)}{1+t^2}.
\]
The factor $1+t^2$ cancels the factor in $dx$. The hypothesis $a>|b|$ keeps the denominator from vanishing. At $x=\pi$ the half-angle tangent is unbounded, so the definite integral is the improper integral in $t$.

\texttt{U26-D03} is $a=3$, $b=2$, stored value $\pi/\sqrt{5}$. The feed prints $\sqrt{5}\,\pi/5$. \texttt{U26-D05} is $a=5$, $b=4$, stored value $\pi/3$.

\section{Both sine and cosine --- \texttt{U26-D06}}
\[
\int_0^{\pi/2}\frac{dx}{1+\sin x+\cos x}.
\]
Here $t$ runs from 0 to 1. The numerator $2\,dt/(1+t^2)$ cancels against
\[
1+\sin x+\cos x=\frac{2(1+t)}{1+t^2},
\]
and the integral is $\int_0^1 dt/(1+t)=\ln 2$. That is the stored answer.

\section{Refusal --- a shorter interval, or an antiderivative}
\texttt{U26-D07} is the integrand of \texttt{U26-D02} on $[0,\pi/3]$, not on $[0,\pi]$. Refuse $\pi/\sqrt{3}$. The upper limit in $t$ is $\tan(\pi/6)=1/\sqrt{3}$, which is finite, so the improper integral on $[0,\infty)$ does not apply. The stored value is
\[
\frac{2\sqrt{3}}{3}\arctan\frac{1}{3}.
\]

\texttt{U26-D04} is the indefinite integral of the integrand in \texttt{U26-D01}. Refuse the definite value $1$. The stored antiderivative is $\tan x-\sec x$. On $(0,\pi)$ the half-angle form differs from that expression by the constant $1$:
\[
\tan x-\sec x=-\frac{2}{1+\tan(x/2)}+1.
\]
The drill stores $\tan x-\sec x$.

A pure sine power is unit 20. \texttt{U20-D01} stores $3\pi/16$. A factor $\sec^2 x$ that is already a derivative is \texttt{U29-D08}, stored $1$. Refuse $t=\tan(x/2)$ on those drills.

\section{The early pass}
Unit 09 works one integral before this drill set. On $[0,\pi/3]$,
\[
\int\frac{dx}{1+\cos x}=\tan\frac x2,
\]
and the definite value is $\sqrt{3}/3$. That value is not one of \texttt{U26-D01} through \texttt{U26-D08}. The same cancellation, $1+\cos x=2/(1+t^2)$, is the pattern to recognise. A miss of the substitution still teaches this file, then returns to that early integral if the limits were the part that failed. The early-pass blank writes the limits of \texttt{U26-D01}. The faded dose inside this file is \texttt{U26-D08}.

\section{What this is not}
A reduction $I_n=((n-1)/n)I_{n-2}$ is unit 20. Adding a reflected copy is unit 10. Introducing a constant in a rational function of $x$ is \texttt{docs/techniques/introduce-constant.md}. If the integrand is not rational in sine and cosine, this substitution makes a harder integral. A rational function of $\tan x$ that already contains $\sec^2 x$ is $u=\tan x$, unit 09, not $t=\tan(x/2)$.

\section{Dose}
Worked reads \texttt{U26-D01}, then the faded limits on \texttt{U26-D08}, then sets \texttt{U26-D02}. Faded gives only those two limits. Bare sets \texttt{U26-D02} and hides $a$ and $b$. Quoting the $[0,\pi]$ value on a shorter interval opens the refusal and sets \texttt{U26-D07}. A definite number written as an antiderivative opens the refusal and sets \texttt{U26-D04}.

\section{Drill pointers}
See \texttt{drills.tex} / \texttt{drills.md}. Feed: \texttt{verify/drills\_unit26.json}.

\section*{On a miss}
Teach this file. Write the three replacements, change the limits, and simplify before you expand. If the interval is shorter than $[0,\pi]$, set \texttt{U26-D07}. If the blank wants $+C$, differentiate \texttt{U26-D04}. Do not invent an answer that is not in \texttt{drills.md}.
\end{document}
